\section*{Hoofdstuk \ref{ch:b1:lab-techniques-1} --- Labotechnieken I: veiligheid, meten en scheiden}

\begin{solution}{exo:b1:lab-techniques-1:1}
\emph{Gevaar}: H225 (licht ontvlambare vloeistof, categorie 2) vereist
het, en het etiket draagt het ernstigste woord. H225 is een fysisch
\omterm{def:b1:lab-techniques-1:hazard}{gevaar}; H319 (ernstige oogirritatie) en H336 (slaperigheid of
duizeligheid) zijn \omterm{def:b1:lab-techniques-1:hazard}{gevaren} voor de gezondheid. Voorzorgen: geen vlam of
vonk in de buurt, een gesloten fles; een veiligheidsbril; werken in een
geventileerde ruimte of een zuurkast.
% ledger: ghs:acetone, hstat:H225, hstat:H319, hstat:H336
\end{solution}

\begin{solution}{exo:b1:lab-techniques-1:2}
(a) Een \omterm{def:b1:lab-techniques-1:hazard}{gevaar}, een eigenschap van het zuur. (b) Een \omterm{def:b1:lab-techniques-1:hazard}{risico}, klein omdat
de blootstelling klein is. (c) Een \omterm{def:b1:lab-techniques-1:hazard}{gevaar}. (d) Een \omterm{def:b1:lab-techniques-1:hazard}{risico}: hetzelfde
\omterm{def:b1:lab-techniques-1:hazard}{gevaar}, een grotere blootstelling (damp, spatten) als het heet en open is.
\end{solution}

\begin{solution}{exo:b1:lab-techniques-1:3}
$R_f = 2.1/6.0 = 0.35$ en $4.2/6.0 = 0.70$. De tweede vlek heeft dezelfde
$R_f$ als de referentie: het monster kan die stof bevatten (het is
consistent), maar gelijke $R_f$-waarden met één eluens bewijzen de
identiteit niet; een tweede eluens, of een gemengde vlek van monster en
referentie die één enkele vlek blijft, versterkt het besluit. De vlek op
0.35 is zeker een ander deeltje.
\end{solution}

\begin{solution}{exo:b1:lab-techniques-1:4}
Halve breedte \qty{0.05}{mL}, rechthoekig: $u = 0.05/\sqrt3 = \qty{0.029}{mL}$.
Voor een verschil van twee onafhankelijke aflezingen is
$u = \sqrt{2} \times 0.029 = \qty{0.041}{mL}$.
\end{solution}

\begin{solution}{exo:b1:lab-techniques-1:5}
Gemiddelde \qty{0.25100}{g}; afwijkingen $+2$, $-2$, $+5$, 0, $-5$ (in
\qty{e-4}{g}), som van de kwadraten \qty{58e-8}{g^2}, $s = \sqrt{58 \times
10^{-8}/4} = \qty{3.8e-4}{g}$; $u = s/\sqrt5 = \qty{1.7e-4}{g}$;
$U = \qty{3.4e-4}{g}$: $m = (0.25100 \pm 0.00034)$~g, $k = 2$.
\end{solution}

\begin{solution}{exo:b1:lab-techniques-1:6}
Eén portie: $f = 1/(1 + 3 \times 30/50) = 0.357$, \qty{64.3}{\percent}
geëxtraheerd. Twee van \qty{15}{mL}: $f = 1.9^{-2} = 0.277$,
\qty{72.3}{\percent}. Drie van \qty{10}{mL}: $f = 1.6^{-3} = 0.244$,
\qty{75.6}{\percent}.
\end{solution}

\begin{solution}{exo:b1:lab-techniques-1:7}
$z = 0.0012/\sqrt{0.0006^2 + 0.0002^2} = 0.0012/0.00063 = 1.9 \le 2$:
compatibel, net. Met \qty{0.0004}{mol/L}: $z = 0.0012/0.00045 = 2.7$,
niet compatibel: de preciezere \omterm{def:b1:titration-methods:titration}{titratie} brengt een echt verschil aan het
licht (een fout bij de bereiding, of bij de \omterm{def:b1:titration-methods:titration}{titratie}).
\end{solution}

\begin{solution}{exo:b1:lab-techniques-1:8}
Q: koud weinig oplosbaar, warm zeer goed oplosbaar. P lost de vaste stof
warm nauwelijks beter op dan koud (er kristalliseert niets bij het
afkoelen); R houdt koud te veel opgelost (grote verliezen). Met Q vraagt
\qty{5.0}{g} minstens $5.0/12 \times 100 = \qty{42}{mL}$ kokend
oplosmiddel, dat koud $0.3 \times 0.417 = \qty{0.13}{g}$ opgelost houdt:
hoogstens \qty{4.9}{g}, \qty{97.5}{\percent}, kan teruggewonnen worden
(vóór elk verlies op de filter).
\end{solution}

\begin{solution}{exo:b1:lab-techniques-1:9}
Ethanol ($n_D = 1.3611$ bij \qty{20}{\celsius}); water (1.333) en
cyclohexaan (1.4266) zijn uitgesloten. De index varieert met de
temperatuur, dus een waarde zegt niets zonder haar. Een waarde van 1.350,
tussen water en ethanol, wijst op een mengsel, zoals ethanol dat water
bevat.
% ledger: nD:water, nD:ethanol, nD:cyclohexane
\end{solution}

\begin{solution}{exo:b1:lab-techniques-1:10}
Met $t = KV/V_a$ en $n$ porties is $\ln f_n = -n\ln(1 + t/n)$; voor
$n \to \infty$ gaat $n \ln(1 + t/n) \to t$ (omdat $\ln(1 + \varepsilon) \sim
\varepsilon$), en daalt $f_n$ naar $\mathrm{e}^{-t}$ (\cref{prop:b1:lab-techniques-1:multiple-extractions}).
Hier is $t = 3 \times 30/50 = 1.8$, $\mathrm{e}^{-1.8} = 0.165$: hoogstens
\qty{83.5}{\percent}. Drie porties geven al \qty{75.6}{\percent},
\qty{91}{\percent} van de limiet; \qty{99}{\percent} ervan bereiken vraagt
32 porties van minder dan \qty{1}{mL}: voorbij twee of drie porties is de
winst het werk niet waard.
\end{solution}

\begin{solution}{exo:b1:lab-techniques-1:11}
Met $t = x - x_0$: $\int_{-a}^{a} (a - |t|)/a^2 \,\mathrm{d}t =
2 \int_0^a (a - t)/a^2 \,\mathrm{d}t = 2 (a^2/2)/a^2 = 1$. Het gemiddelde
is $x_0$ wegens de symmetrie; de variantie is $2 \int_0^a t^2 (a - t)/a^2
\,\mathrm{d}t = \frac{2}{a^2}\left(\frac{a^4}{3} - \frac{a^4}{4}\right) =
\frac{a^2}{6}$, dus $u = a/\sqrt6 = 0.41\,a$, tegenover $a/\sqrt3 = 0.58\,a$
in het rechthoekige geval: weten dat het midden waarschijnlijker is,
verkleint de onzekerheid.
\end{solution}

\begin{solution}{exo:b1:lab-techniques-1:12}
$c = m/(MV) = 0.25100/(204.2 \times 0.10000) = \qty{0.012292}{mol/L}$.
Relatieve onzekerheden: massa $1.7 \times 10^{-4}/0.2510 = 6.8 \times
10^{-4}$, volume $0.06/100.00 = 6.0 \times 10^{-4}$; gecombineerd
$\sqrt{(6.8^2 + 6.0^2)} \times 10^{-4} = 9.1 \times 10^{-4}$, dus
\qty{0.091}{\percent}; $u(c) = \qty{1.1e-5}{mol/L}$, $U = \qty{2.2e-5}{mol/L}$:
$c = (0.012292 \pm 0.000022)$~mol/L, $k = 2$. De weging overheerst,
licht; de twee bronnen zijn van dezelfde grootte.
\end{solution}

\begin{solution}{pb:b1:lab-techniques-1:1}
\textbf{1.} GHS02 ontvlambaar, GHS05 bijtend, GHS06 acute toxiciteit,
GHS07 schadelijk of irriterend.
\textbf{2.} Salicylzuur: \emph{Gevaar}, wegens H318 (ernstig oogletsel,
categorie 1). Aspirine: \emph{Waarschuwing} (alleen H302).
\textbf{3.} Het \omterm{def:b1:lab-techniques-1:hazard}{gevaar} ligt vast, de blootstelling niet: een klein
volume, afgemeten en gebruikt in een zuurkast, met handschoenen,
veiligheidsbril en labojas, de fles meteen gesloten, geen vlam in de
buurt.
\textbf{4.} H314 (ernstige brandwonden en oogletsel): veiligheidsbril en
handschoenen. H332 (schadelijk bij inademing) en H226 (ontvlambare
vloeistof en damp): de zuurkast en de afwezigheid van vlammen.
\textbf{5.} Voorzichtig gedurende een aantal minuten met water afspoelen,
contactlenzen verwijderen indien mogelijk, en blijven spoelen
(P305+P351+P338); daarna een arts raadplegen.
\textbf{6.} In het waterige afval, na neutralisatie, niet door de
gootsteen.
\textbf{7.} \ce{C7H6O3 + C4H6O3 -> C9H8O4 + C2H4O2}: C $7 + 4 = 9 + 2$,
H $6 + 6 = 8 + 4$, O $3 + 3 = 4 + 2$.
\textbf{8.} $2.00/138.0 = \qty{1.449e-2}{mol}$; hoogstens
$\num{1.449e-2} \times 180.0 = \qty{2.61}{g}$ aspirine.
\textbf{9.} De ruwe vaste stof bevat niet-gereageerd salicylzuur,
ethaanzuur en water: haar massa is niet die van aspirine, en haar
zuiverheid is onbekend.
\textbf{10.} Wat na het afkoelen opgelost blijft, gaat verloren: een
\omterm{def:b1:precipitation:solubility}{oplosbaarheid} die koud klein en warm groot is, betekent dat het product
warm oplost en koud bijna volledig terugkomt.
\textbf{11.} Elke extra milliliter houdt koud zijn deel van het product
opgelost.
\textbf{12.} Langzame groei geeft grotere, regelmatigere \omterm{def:b1:crystals-metals:crystal}{kristallen} die
minder moederloog en onzuiverheden insluiten; het ijs verlaagt de
\omterm{def:b1:precipitation:solubility}{oplosbaarheid}, en dus het verlies.
\textbf{13.} $\qty{0.045}{L} \times \qty{3.3}{g/L} = \qty{0.15}{g}$.
\textbf{14.} $1.95/2.61 = \qty{75}{\percent}$ (\qty{74.7}{\percent}).
\textbf{15.} $807/6 = \qty{134.5}{\celsius}$.
\textbf{16.} Afwijkingen $-0.5$, $0.5$, $-1.5$, $0.5$, $-0.5$, $1.5$; som
van de kwadraten 5.5; $s = \sqrt{5.5/5} = \qty{1.05}{\celsius}$.
\textbf{17.} $u_A = 1.05/\sqrt6 = \qty{0.43}{\celsius}$.
\textbf{18.} Halve breedte \qty{0.5}{\celsius}: $0.5/\sqrt3 = \qty{0.29}{\celsius}$.
\textbf{19.} $1/\sqrt3 = \qty{0.58}{\celsius}$;
$u = \sqrt{0.428^2 + 0.289^2 + 0.577^2} = \sqrt{0.600} = \qty{0.77}{\celsius}$.
\textbf{20.} $U = 2 \times 0.775 = \qty{1.5}{\celsius}$:
$T = (134.5 \pm 1.5)$~\unit{\celsius}, $k = 2$.
\textbf{21.} Het smelt laag en over een gebied van \qty{8}{\celsius}: het
is onzuiver (salicylzuur, ethaanzuur, water).
\textbf{22.} Tot op de eenheid gegeven ligt de waarde binnen
$\pm\qty{0.5}{\celsius}$: $u_{\mathrm{ref}} = 0.5/\sqrt3 = \qty{0.29}{\celsius}$.
\textbf{23.} $2.3/6.0 = 0.38$ en $3.3/6.0 = 0.55$. Het ruwe product bevatte
salicylzuur (dezelfde $R_f$ als de referentie) en aspirine; na de
\omterm{def:b1:lab-techniques-1:recrystallisation}{herkristallisatie} blijft alleen de vlek van aspirine over: het
salicylzuur is verwijderd, althans tot onder wat de plaat kan detecteren.
\textbf{24.} Aspirine ontleedt bij het smelten; langzaam verwarmd heeft
ze tijd om te ontleden, en het gevormde mengsel smelt lager. De
getabelleerde waarde slaat op snelle verwarming, zoals op de bank.
\textbf{25.} $z = |134.5 - 135|/\sqrt{0.600 + 0.083} = 0.5/0.827 =$
\textbf{0.60}: ruim onder 2, het gemeten en het getabelleerde smeltpunt
zijn compatibel.
\end{solution}
